\chapter*{Prefaci}
\addcontentsline{toc}{chapter}{Prefaci}

Aquest llibre forma part d'\textbf{aprendes}, una col·lecció de materials
d'estudi que comparteixen una mateixa manera de treballar: cada tema es
presenta en una part de \emph{teoria}, es reprèn en una part de
\emph{pràctica} amb problemes resolts, es consolida amb \emph{exercicis
proposats} i es comprova amb \emph{\nomtests}. Qui ja conegui
altres volums de la col·lecció hi reconeixerà l'estructura; qui s'hi
acosti per primera vegada la trobarà explicada a la guia de lectura.

El volum té un propòsit precís. És el primer d'un projecte en dues
etapes: aquí es construeix el llenguatge de la teoria de categories
---des de les definicions bàsiques fins als classificadors de subobjectes
i els topos--- amb el detall que necessitarà un volum posterior dedicat
a la lògica categòrica. Per això la lògica hi és present des del
començament, com a font d'exemples i com a direcció del recorregut,
tot i que el seu desenvolupament sistemàtic queda per a aquella segona
etapa.

\ifpublicacioparcial
Aquesta és una edició parcial, que conté els tres primers temes
---categories, functors i transformacions naturals--- en les quatre parts,
juntament amb els apèndixs i els materials de consulta. Els capítols
restants s'incorporaran en l'edició completa.
\else
Aquesta edició conté el volum complet: nou temes desenvolupats en les
quatre parts i un capítol final que fa de pont cap a la lògica categòrica.
\fi

Abans del primer capítol, la \emph{Introducció} explica per què val la
pena aprendre aquest llenguatge i la \emph{Guia de lectura} proposa com
fer-ho. Es poden llegir en aquest ordre o tornar-hi quan convingui.

\medskip
\noindent\textbf{Prerequisits.} El text pressuposa la maduresa matemàtica habitual d'un primer curs universitari de matemàtiques, informàtica o lògica: treballar amb conjunts i aplicacions, seguir una demostració elemental i llegir notació algebraica bàsica. No es pressuposa haver estudiat teoria de categories, topologia, àlgebra abstracta avançada ni lògica categòrica. Quan un exemple necessita vocabulari addicional, s'introdueix localment o es remet a un apèndix.

\medskip
\noindent\textbf{Errates i suggeriments.} Són benvinguts. Es poden comunicar obrint una incidència al repositori del projecte, \llibreerrates, indicant la versió (\llibreversio) i la pàgina.

% ------------------------------------------------------------
% Espai reservat per a l'autor (descomenteu i completeu):
%
% \medskip
% \noindent\textbf{Agraïments.} ...
%
% (Les errates i els suggeriments: vegeu la pàgina de crèdits i el text següent.)
% ------------------------------------------------------------
